% !Mode:: "TeX:UTF-8"
\fancyhead[CO]{\song\wuhao ~\chaptername \quad 总结与展望~}
\fancypagestyle{plain}{\fancyhead[CO]{\song\wuhao ~\chaptername \quad 总结与展望~}}

\chapter{总结与展望}

% 第一节
\section{总结}

Transformer 模型在自然语言处理、大模型推理中的使用范围已经扩大，让矩阵乘法和多头注意力等主要算子的实现效率变得重要，这会影响到模型部署的性能。异构硬件平台的发展加快了速度，传统依靠人工优化内核和厂商专用库的实现方式不得不承担较高的开发与移植，并且维护成本也变得比较大。现在深度学习编译框架在块级逻辑表达、跨平台优化、性能迁移方面有不足。本文运用 MLIR，解决 Transformer 模型主要算子的自动代码生成问题，并且设计及实现了面向异构平台的高性能算子生成框架 DeepGen。

本文的主要内容主要包括下面这些方面。本文提出面向 Transformer 主要算子的块级分层中间表示方式，以解决现有编译框架在生成高性能算子时表达能力不足的问题。本文凭借 MLIR 的多级表示机制设计了 DeepGen Dialect 用来对矩阵乘法、注意力机制里的块级计算逻辑、数据搬运行为和并行结构进行显式描述，这给后续优化和降级提供了统一的表示基础。

针对块级中间表示向线程级实现映射的问题，本文设计了依靠降级信息（LowerInfo）进行推断的转换方式。该方式对缓冲区布局、线程映射关系和索引规则的结构进行推导，以实现从高层块级描述到底层线程执行逻辑的渐进式转换，并且可以支持访存优化、并行调度和代码生成。

针对异构硬件平台在执行模型和存储层次上的差异，本文提出了一套通用的跨平台优化方案，包括 Software Pipeline、Local Split K、Block Swizzle 和 Access Scatter 等技术。这些方法通过计算调度、规约拆分、线程块映射和共享存储访问等操作，使生成软件代码的执行效率得以提高，同时也改善了该框架对各类适配硬件平台的能力。

针对调优空间较大和搜索成本较高的问题，本文设计了通过空间剪枝实现自动调优的方式。该方式依靠线程块、线程负载分块等参数来构成搜索空间，并且利用性能测试缓存和参数遍历策略，自动得到较优的配置，这样可以减少人工进行参数调整的工作量。

实验结果说明，DeepGen 在 NVIDIA GPU 和海光 DCU 平台上都拥有较好的性能和可移植性。在矩阵乘法测试中，相比 cuBLAS 和 rocBLAS，本文实现平均 1.04 倍和 1.03 倍性能提高。在注意力机制测试中，相较于 Triton 在 NVIDIA 和海光平台上平均加速分别可以达到 2.55 倍和 3.09 倍。在端到端推理中，本文对 BERT、GPT2 和 LLaMA2 等模型进行实验，同样获得了稳定的收益。精度测试的结果表明，生成算子的输出结果和 PyTorch 基准实现保持了高度一致性。

此外，自动调优实验结果表明，本文提出的调优空间剪枝策略能够有效减少无效候选配置数量，降低自动调优过程中的搜索开销。优化 Pass 消融实验进一步说明，不同优化 Pass 的收益具有明显的平台相关性和配置相关性，不存在对所有场景均最优的固定优化组合，因此将优化 Pass 纳入自动调优空间进行自动选择具有必要性。

% 第二节
\section{展望}

尽管本文针对 Transformer 模型中主要算子的自动代码生成问题，设计并实现了使用 MLIR 的 DeepGen 框架，并且达到了比较理想的实验效果，但是仍然存在进一步优化的余地。

在低精度混合精度代码生成方面，DeepGen 仅支持 FP32 精度，尚未充分覆盖 FP16、BF16 等低精度场景。在大模型推理部署中，低精度计算通常能够显著提高矩阵乘法和注意力机制等核心算子的吞吐率，但其高性能实现往往依赖硬件专用矩阵计算单元以及相应的数据布局和指令映射策略。后续工作将进一步扩展 DeepGen 的低精度代码生成能力，在现有双层 IR 和优化 Pass 的基础上，引入面向 Tensor Core、Matrix Core 等矩阵计算单元的后端映射机制，从而支持 FP16/BF16 等混合精度下的计算。

在新型硬件单元适配方面，DeepGen 目前仅支持 CUDA Core 执行路径，对新型异步 Tensor Core 和 TMA 等异步硬件机制的支持还不够充分。由于硬件架构在不断演进，异步执行已成为提高计算和数据搬运重叠度的主要方式。后续工作可以在现有分层中间表示的基础上，增强对异步拷贝、异步矩阵乘及多阶段流水执行的表示能力，并且设计相应的降低策略与优化 Pass，从而提高该框架在混合精度和异步执行场景下的性能上限。

在自动调优方面，目前 DeepGen 主要还是使用参数遍历和规则剪枝方法。参数空间如果一直扩大，硬件特征变得更加复杂，该方式在效率和搜索质量上会表现出不足。后续本文可以将启发式搜索、历史结果驱动的增量搜索、性能预测模型等方法进行引入，以便减少无效编译和测试的次数，从而提高调优效率。

未来还可以把优化范围从算子级扩展到图级，分析跨算子融合、运行时调度共同优化和整的所有图推理优化，让代码生成能力、图级编译优化之间进行结合变得更加紧密。


